\documentclass[11pt]{article}

\usepackage[a4paper,margin=1in]{geometry}
\usepackage{amsmath,amssymb,mathtools}
\usepackage[hidelinks]{hyperref}

\title{Level-Set Deficit Identity for the Nicola--Tilli Route}
\author{}
\date{}

\begin{document}
\maketitle

This note implements the first concrete step from \texttt{plan2.md}: turn the
convexity gap in \texttt{faber-krahn-1.pdf} into an exact weighted integral of
the two losses in the level-set proof.

Throughout, let \(u=u_\Omega\) solve
\[
-\Delta u=1 \quad\hbox{in }\Omega,\qquad
u=0\quad\hbox{on }\partial\Omega,
\]
and set
\[
E_t=\{u>t\},\qquad m(t)=|E_t|,\qquad
\Sigma_t=\{u=t\},\qquad P(t)=\mathcal H^{n-1}(\Sigma_t).
\]
The formulas below are written for smooth domains and regular levels. The same
identities should extend to the finite-perimeter/coarea setting by the usual
approximation and a.e. level-set interpretation.

\section{The Two Pointwise Losses}

For a regular level \(t\), define
\[
a(t):=-m'(t)=\int_{\Sigma_t}\frac{1}{|\nabla u|}\,d\mathcal H^{n-1}.
\]
By the divergence theorem applied to \(E_t\),
\[
\int_{\Sigma_t}|\nabla u|\,d\mathcal H^{n-1}=m(t).
\]
Let
\[
c_n:=n^2\omega_n^{2/n}.
\]
The total level-set loss in the differential inequality is
\[
D_{\rm lev}(t):=a(t)m(t)-c_n m(t)^{2-2/n}.
\]
It splits as
\[
D_{\rm lev}(t)=D_H(t)+D_I(t),
\]
where
\[
D_H(t)
:=
\left(\int_{\Sigma_t}\frac{1}{|\nabla u|}\right)
\left(\int_{\Sigma_t}|\nabla u|\right)-P(t)^2
=a(t)m(t)-P(t)^2
\]
is the Holder/Cauchy--Schwarz loss, and
\[
D_I(t):=P(t)^2-c_n m(t)^{2-2/n}
\]
is the squared-perimeter isoperimetric loss of \(E_t\).

Both terms are nonnegative: \(D_H(t)\ge0\) by Cauchy--Schwarz, and
\(D_I(t)\ge0\) by the isoperimetric inequality.

\section{Convexity Variable and Exact Identity}

Let \(L:=|\Omega|\), let \(v(s)=u^*(s)=m^{-1}(s)\), and define
\[
I(s):=\int_{\{u>v(s)\}}u\,dx.
\]
Equivalently, for regular values,
\[
I'(s)=v(s),\qquad
I''(s)=v'(s)=\frac{1}{m'(v(s))}=-\frac{1}{a(v(s))}.
\]
The function used in \texttt{faber-krahn-1.pdf} is
\[
G(s):=
I(s)+\frac{s^{1+2/n}}{(4+2n)\omega_n^{2/n}}.
\]
Since
\[
\frac{d^2}{ds^2}
\frac{s^{1+2/n}}{(4+2n)\omega_n^{2/n}}
=
\frac{1}{c_n s^{1-2/n}},
\]
we get, with \(s=m(t)\),
\[
G''(s)
=
\frac{1}{c_n s^{1-2/n}}-\frac1{a(t)}
=
\frac{D_{\rm lev}(t)}
{s\,a(t)\,c_n s^{1-2/n}}.
\]
The endpoint convexity gap is
\[
\Gamma(\Omega):=G'(L)-\frac{G(L)}{L}.
\]
Because \(G(0)=0\),
\[
\Gamma(\Omega)=\frac1L\int_0^L sG''(s)\,ds.
\]
Changing variables \(s=m(t)\), so that \(ds=-a(t)\,dt\), gives the exact
deficit identity
\[
\boxed{
\Gamma(\Omega)
=
\frac1{|\Omega|}
\int_0^{\|u\|_\infty}
\frac{D_H(t)+D_I(t)}
{n^2\omega_n^{2/n}m(t)^{1-2/n}}\,dt.
}
\]
This is the clean quantitative form of the rigidity statement at the end of
\texttt{faber-krahn-1.pdf}.

\section{Relation to the Saint-Venant Deficit}

At \(s=L\), one has \(v(L)=0\). Hence
\[
G'(L)=\frac{L^{2/n}}{2n\omega_n^{2/n}}.
\]
Also
\[
\frac{G(L)}{L}
=
\frac1L\int_\Omega u\,dx
+
\frac{L^{2/n}}{(4+2n)\omega_n^{2/n}}.
\]
Therefore
\[
\Gamma(\Omega)
=
\frac1L\left(\int_B u_B\,dx-\int_\Omega u_\Omega\,dx\right),
\]
where \(B\) is the ball with \(|B|=|\Omega|\). Since
\[
E(\Omega)=-\frac12\int_\Omega u_\Omega\,dx,
\]
we obtain
\[
\boxed{
\Gamma(\Omega)
=
\frac{2}{|\Omega|}
\big(E(\Omega)-E(B)\big).
}
\]
Thus the endpoint convexity gap is not merely comparable to the torsion
deficit; it is exactly the normalized Saint-Venant deficit.

\section{Profile-Gap Monotonicity Inspired by the STFT Stability Paper}

The paper \texttt{s00222-024-01248-2-3.pdf} proves stability of the
Nicola--Tilli STFT theorem by studying the area between the rearranged profile
\(u^*(s)\) and the optimizer profile \(e^{-s}\). There is a direct torsion
analogue.

Let \(B\) be the ball with \(|B|=|\Omega|=:L\), and let
\[
b(s):=u_B^*(s),\qquad v(s):=u_\Omega^*(s).
\]
For a ball of radius \(R\), \(L=\omega_n R^n\), one has
\[
b(s)=\frac{R^2-(s/\omega_n)^{2/n}}{2n},
\qquad 0\le s\le L,
\]
and therefore
\[
b'(s)=-\frac{1}{n^2\omega_n^{2/n}s^{1-2/n}}.
\]
The level-set differential inequality gives
\[
v'(s)\ge b'(s)
\]
for a.e. \(s\in(0,L)\). Hence
\[
h(s):=b(s)-v(s)
\]
is nonnegative and nonincreasing, with \(h(L)=0\). Moreover
\[
\int_0^L h(s)\,ds
=
\int_B u_B\,dx-\int_\Omega u_\Omega\,dx
=
2(E(\Omega)-E(B)).
\]
Consequently, for every \(s\in(0,L]\),
\[
0\le b(s)-v(s)
\le
\frac{2(E(\Omega)-E(B))}{s}.
\]
In particular, on every boundary-volume slab \(s\in[\theta L,L]\),
\[
\|b-v\|_{L^\infty([\theta L,L])}
\le
\frac{2(E(\Omega)-E(B))}{\theta L}.
\]
This is a useful quantitative replacement for compactness at the one-dimensional
profile level. It says that small Saint-Venant deficit forces the torsion level
corresponding to a prescribed near-boundary volume \(s\sim L\) to be close to the
ball's corresponding level.

For example, if \(s=L-\eta\) and \(\eta\ll L\), then
\[
b(L-\eta)
=
\frac{R^2-((L-\eta)/\omega_n)^{2/n}}{2n}
\simeq_{n,L}\eta.
\]
Thus if \(E(\Omega)-E(B)\ll \eta\), the level
\[
t_\eta:=v(L-\eta)
\]
is positive and comparable to the ball boundary-layer height. This provides a
candidate near-boundary level \(E_{t_\eta}\) with exactly
\[
|\Omega\setminus E_{t_\eta}|=\eta.
\]
The missing part is still to make sure this same level has small geometric
level-set deficit. The exact identity above gives averaged control, and the
hybrid strategy is to combine this profile control with selected-minimizer
regularity to extract good near-boundary levels.

\section{Consequences for Good Levels}

The quantitative isoperimetric inequality gives, for a dimensional constant
\(c>0\),
\[
D_I(t)
\ge
c\,m(t)^{2-2/n}\mathcal A(E_t)^2
\]
for a.e. regular level \(t\). Hence
\[
E(\Omega)-E(B)
\gtrsim_n
\int_0^{\|u\|_\infty}
m(t)\mathcal A(E_t)^2\,dt
\]
up to the normalization factor \(|\Omega|\). More precisely, the exact identity
controls the weighted average
\[
\int_0^{\|u\|_\infty}
\frac{D_I(t)}
{m(t)^{1-2/n}}\,dt.
\]

For any measurable level interval \(J\subset(0,\|u\|_\infty)\), define the
weight
\[
d\nu(t)
:=
\frac{dt}{n^2\omega_n^{2/n}m(t)^{1-2/n}}.
\]
If \(\nu(J)>0\), then there exists a regular \(t\in J\) such that
\[
D_H(t)+D_I(t)
\le
\frac{|\Omega|\Gamma(\Omega)}{\nu(J)}.
\]
In particular, on any slab where \(m(t)\) is bounded above and below, the
Saint-Venant deficit produces at least one level with both small isoperimetric
loss and small Holder loss.

\section{Holder Loss as a Serrin-Variance Defect}

Let \(f=|\nabla u|\) on \(\Sigma_t\), and set
\[
\bar f:=\frac1{P(t)}\int_{\Sigma_t}f
=
\frac{m(t)}{P(t)}.
\]
Then the Holder deficit has the exact weighted variance representation
\[
\boxed{
\int_{\Sigma_t}
\frac{(f-\bar f)^2}{f}\,d\mathcal H^{n-1}
=
\frac{m(t)}{P(t)^2}D_H(t).
}
\]
Indeed,
\[
\int_{\Sigma_t}\frac{(f-\bar f)^2}{f}
=
\int_{\Sigma_t}f
-2\bar f P(t)
+\bar f^2\int_{\Sigma_t}\frac1f
=
\frac{m(t)}{P(t)^2}
\left(
m(t)\int_{\Sigma_t}\frac1f-P(t)^2
\right).
\]
Thus small \(D_H(t)\) says that \(|\nabla u|\) is almost constant on
\(\Sigma_t\), in the natural weighted \(L^2(1/f)\) sense.

If, on a good regular level, one also has
\[
0<c_t\le |\nabla u|\le C_t<\infty,
\]
then this becomes an ordinary \(L^2\)-variance estimate:
\[
\int_{\Sigma_t}|f-\bar f|^2\,d\mathcal H^{n-1}
\le
C_t\frac{m(t)}{P(t)^2}D_H(t).
\]
This is the precise bridge from the Nicola--Tilli convexity gap to an
approximate Serrin condition on the interior domain \(E_t\).

\section{What Remains Open in This Route}

The identity proves that small Saint-Venant deficit forces many superlevel sets
to be simultaneously almost isoperimetric and almost overdetermined.

The unsolved propagation problem is to recover \(\mathcal A(\Omega)^2\) from
this information. The deficit identity sees interior levels with positive
height, while Fraenkel asymmetry is a \(t=0\) boundary datum. A sharp proof by
this route needs a boundary-layer theorem of the following form:
\[
\exists t\downarrow0
\quad\hbox{such that}\quad
|\Omega\setminus E_t|
\hbox{ is small and }
D_H(t)+D_I(t)
\hbox{ remains controlled.}
\]
Without such a theorem, the level-set method gives strong information on
interior cores \(E_t\), but not yet a sharp estimate for the original domain
\(\Omega\).

\end{document}
